\section{Fusion isometries}

The next definitions isolate the transfer-map content underlying the
fusion-isometry picture of \cite[Section~4.5]{Cirac2016MPDO_arXiv}: they record
when the blocked transfer map factors through its support algebra as a retract,
which is the idempotence criterion, not the physical fusion isometry of two
tensors into one. For each blocked size $n \ge 1$, the doubled-index blocked
tensor of an MPO has a blocked transfer map $E_n$ acting on bond-space matrices,
and the support algebra is modeled by a subspace through which $E_n$ factors as
a retract. The physical fusion isometries of
\cite[Theorem~4.14(iii)]{Cirac2016MPDO_arXiv} are recorded at the end of this
section, together with the derivation of the same-length product law from
them.

\begin{definition}[Transfer-retract datum]%
    \label{def:mpdo_fusion_isometry}
    \lean{MPOTensor.TransferRetractData}
    \leanok
    \uses{def:mpo_tensor}
    A \emph{transfer-retract datum} at blocked size $n$ for an
    MPO tensor consists of a subspace $\mathcal{A}_n \subseteq \MN{D}$
    together with linear maps
    \begin{align}
        T_n&:\MN{D}\longrightarrow\mathcal{A}_n,
        &
        S_n&:\mathcal{A}_n\longrightarrow\MN{D},
        \notag\\
        T_n\circ S_n&=\id_{\mathcal{A}_n},
        &
        S_n\circ T_n&=E_n,
        \label{eq:rfp_transfer_retract}
    \end{align}
    where $E_n$ denotes the blocked transfer map.
\end{definition}

\begin{definition}[Transfer-retract RFP predicate]%
    \label{def:mpdo_rfp_via_fusion}
    \lean{MPOTensor.IsRFP_MPDO_via_transferRetract}
    \leanok
    \uses{def:mpdo_fusion_isometry, def:mpo_tensor}
    An MPO tensor $M$ satisfies the \emph{transfer-retract formulation} of the
    renormalization fixed-point condition when for every blocked size $n \ge 1$
    there exists transfer-retract data for the blocked transfer
    map $E_n$.
\end{definition}

\begin{theorem}[Transfer-retract data imply transfer-map idempotence]%
    \label{thm:mpdo_rfp_of_fusion}
    \lean{MPOTensor.TransferRetractData.isRFP, MPOTensor.isRFP_of_isRFP_MPDO_via_transferRetract}
    \leanok
    \uses{def:mpdo_rfp_via_fusion}
    A one-site transfer-retract datum implies transfer-map idempotence. Hence the
    transfer-retract formulation, which supplies such data at every positive
    blocked size, implies the same condition.
\end{theorem}
\begin{proof}\leanok
    Apply the definition at blocked size $n=1$.
    By (\ref{eq:rfp_transfer_retract}),
    \begin{align}
        E_1^2
        &=S_1T_1S_1T_1
          =S_1(T_1S_1)T_1
          =S_1T_1
          =E_1.
        \notag
    \end{align}
    Since $E_1$ is the original transfer map, this is exactly transfer-map
    idempotence.
\end{proof}

\begin{theorem}[One-site transfer-retract criterion]%
    \label{thm:mpdo_one_site_fusion_iff_rfp}
    \lean{MPOTensor.transferRetractData_one_iff_isRFP}
    \leanok
    \uses{def:mpdo_fusion_isometry}
    A one-site transfer-retract datum exists if and only if the MPO
    transfer map is idempotent.
\end{theorem}
\begin{proof}\leanok
    The forward implication is the retract calculation
    $E_1^2=S_1T_1S_1T_1=S_1T_1$. Conversely, if $E_1$ is idempotent, then
    it factors through its range, giving the required one-site datum.
\end{proof}

\begin{theorem}[Transfer-retract criterion for idempotence]%
    \label{thm:mpdo_rfp_via_fusion_iff}
    \lean{MPOTensor.isRFP_MPDO_via_transferRetract_iff_isRFP}
    \leanok
    \uses{def:mpdo_rfp_via_fusion}
    The transfer-retract formulation is equivalent to transfer-map idempotence.
\end{theorem}
\begin{proof}\leanok\uses{thm:mpdo_rfp_of_fusion}
    The forward implication is Theorem~\ref{thm:mpdo_rfp_of_fusion}.
    Conversely, if the transfer map $E$ is idempotent, then every positive blocked
    transfer map equals $E$ and therefore factors through its range.
    This range-factorization gives the required transfer-retract data.
\end{proof}

\begin{corollary}[Pure-state recovery]%
    \label{cor:mpdo_fusion_pure_recovery}
    \lean{MPSTensor.toMPOTensor_isRFP_MPDO_via_transferRetract_iff_isTransferIdempotent}
    \leanok
    \uses{thm:mpdo_rfp_via_fusion_iff}
    For the diagonal MPO associated to a pure MPS tensor, the transfer-retract
    formulation reduces to the usual pure-state RFP condition.
\end{corollary}
\begin{proof}\leanok
    Combine Theorem~\ref{thm:mpdo_rfp_via_fusion_iff} with the diagonal pure-state
    recovery theorem for the MPO predicate.
\end{proof}

\begin{theorem}[All-blocked transfer retracts reduce to one site]%
    \label{thm:mpdo_fusion_iff_one_site_fusion}
    \lean{MPOTensor.isRFP_MPDO_via_transferRetract_iff_transferRetractData_one}
    \leanok
    \uses{def:mpdo_rfp_via_fusion, thm:mpdo_one_site_fusion_iff_rfp,
        thm:mpdo_rfp_via_fusion_iff}
    The transfer-retract formulation for all positive blocked sizes is
    equivalent to the existence of a one-site transfer-retract
    datum.
\end{theorem}
\begin{proof}\leanok
    Combine the one-site criterion with the equivalence between the all-blocked
    transfer-retract formulation and transfer-map idempotence.
\end{proof}

\begin{definition}[Operator product of MPO tensors]%
    \label{def:mpdo_operator_product}
    \lean{MPOTensor.mulTensor}
    \leanok
    \uses{def:mpo_tensor}
    The \emph{product tensor} of two MPO tensors contracts the bra index of
    the first factor with the ket index of the second and takes the tensor
    product of the bond spaces:
    \begin{align}
        (M\mathbin{\cdot}N)^{ik}
        &=\sum_j M^{ij}\otimes N^{jk}.
        \label{eq:rfp_product_tensor}
    \end{align}
    In the vertical reading of \cite[Section~4.5]{Cirac2016MPDO_arXiv} this
    is the tensor obtained by joining two vertically read tensors along the
    original horizontal bond, the tensor written $M_\alpha M_\beta$ in
    \cite[Theorem~4.14(iii)]{Cirac2016MPDO_arXiv}.
\end{definition}

\begin{definition}[Canonical reassociation of three bond spaces]%
    \label{def:mpdo_operator_product_associator}
    \lean{MPOTensor.mulTensorAssocEquiv}
    \lean{MPOTensor.mulTensorAssocMatrix}
    \leanok
    \uses{def:mpdo_operator_product}
    For three bond spaces, let
    \begin{align}
        A_{D_1,D_2,D_3}:
        \C^{D_1}\otimes(\C^{D_2}\otimes\C^{D_3})
        &\longrightarrow
        (\C^{D_1}\otimes\C^{D_2})\otimes\C^{D_3}
        \notag
    \end{align}
    be the coordinate pullback induced by canonical reassociation, defined on
    elementary tensors by
    $A_{D_1,D_2,D_3}(x_1\otimes(x_2\otimes x_3))
    =(x_1\otimes x_2)\otimes x_3$.
\end{definition}

\begin{theorem}[Associativity of the product tensor]%
    \label{thm:mpdo_operator_product_associativity}
    \lean{MPOTensor.mulTensor_assoc}
    \lean{MPOTensor.mulTensor_mul_assocMatrix}
    \leanok
    \uses{def:mpdo_operator_product,
        def:mpdo_operator_product_associator}
    The letters of the two parenthesizations of a triple product are related
    by the canonical bond reassociation. Equivalently, for every pair of
    physical indices $i,k$,
    \begin{align}
        ((M\mathbin{\cdot}N)\mathbin{\cdot}P)^{ik}
          A_{D_1,D_2,D_3}
        &=A_{D_1,D_2,D_3}
          (M\mathbin{\cdot}(N\mathbin{\cdot}P))^{ik}.
        \notag
    \end{align}
\end{theorem}
\begin{proof}\leanok
    Expand both product tensors. Reassociation preserves every elementary
    tensor factor, while exchanging the order of the two finite sums gives
    the same contraction over the two intermediate physical indices.
\end{proof}

\begin{lemma}[Word evaluation of the product tensor]%
    \label{lem:mpdo_operator_product_eval_word}
    \lean{MPOTensor.evalWord_mulTensor}
    \leanok
    \uses{def:mpdo_operator_product, def:mpo_eval_word}
    The word evaluation of the product tensor along a configuration pair
    $(\sigma,\tau)$ expands as a sum over the contracted middle
    configurations $\rho$ of tensor products of word evaluations:
    \begin{align}
        (M\mathbin{\cdot}N)^{\sigma\tau}
        &=\sum_\rho M^{\sigma\rho}\otimes N^{\rho\tau}.
        \label{eq:rfp_product_word}
    \end{align}
\end{lemma}
\begin{proof}\leanok
    Induction on the length, using the mixed-product property of the tensor
    product to merge one letter at a time.
\end{proof}

\begin{theorem}[The operator family is multiplicative]%
    \label{thm:mpdo_operator_product_mpo}
    \lean{MPOTensor.mpo_mulTensor}
    \leanok
    \uses{def:mpdo_operator_product, def:mpo_family,
        lem:mpdo_operator_product_eval_word}
    For every chain length $L$, the operator generated by the product tensor
    is the product of the generated operators:
    $O_L(M\mathbin{\cdot}N)=O_L(M)O_L(N)$.
\end{theorem}
\begin{proof}\leanok
    Take traces in (\ref{eq:rfp_product_word}): the trace of a tensor
    product is the product of the traces, and the sum over middle
    configurations is the entrywise form of the operator product.
\end{proof}

\begin{lemma}[Word products of isometrically conjugated letters]%
    \label{lem:mpdo_word_product_conjugation}
    \lean{MPOTensor.listProd_conj_of_conjTranspose_mul_self}
    \leanok
    For an isometry $U$ (a rectangular matrix with $U^\dagger U=1$) and a
    nonempty word of square matrices $F_1,\ldots,F_L$,
    \begin{align}
        (UF_1U^\dagger)(UF_2U^\dagger)\cdots(UF_LU^\dagger)
        &=U F_1F_2\cdots F_L U^\dagger.
        \notag
    \end{align}
\end{lemma}
\begin{proof}\leanok
    The inner factors $U^\dagger U$ telescope to the identity.
\end{proof}

\begin{lemma}[Word products of block-diagonal tensor-product letters]%
    \label{lem:mpdo_word_product_block_kronecker}
    \lean{MPOTensor.listProd_blockDiagonal'_kronecker}
    \leanok
    For block-diagonal letters whose block of index $\gamma$ is the tensor
    product of a fixed matrix $X_\gamma$ with a varying letter
    $G_{\gamma,l}$,
    \begin{align}
        \prod_{l=1}^{L}
            \bigoplus_{\gamma}X_\gamma\otimes G_{\gamma,l}
        &=
        \bigoplus_{\gamma}X_\gamma^{L}\otimes
            \prod_{l=1}^{L}G_{\gamma,l}.
        \notag
    \end{align}
\end{lemma}
\begin{proof}\leanok
    Blocks multiply blockwise and tensor products multiply factorwise, so
    the fixed factors accumulate into the matrix power $X_\gamma^L$.
\end{proof}

\begin{lemma}[Blockwise conjugation and restriction]
    \label{lem:matrix_block_conjugation_restriction}
    \lean{Matrix.blockDiagonal'_conj, Matrix.submatrix_left_conj_equiv}
    \leanok
    Let $F=\bigoplus_a F_a$ and $G=\bigoplus_a G_a$. Then
    $FGF^\dagger=\bigoplus_a F_aG_aF_a^\dagger$.
    Moreover, if $r$ is a bijective reindexing of the row space of a matrix
    $A$, then $A_rXA_r^\dagger=(AXA^\dagger)_{r,r}$, where $A_r$ denotes the
    corresponding row reindexing.
\end{lemma}
\begin{proof}\leanok
    The first identity follows because block-diagonal matrices
    multiply blockwise. The second follows by commuting bijective row
    reindexing with multiplication and adjunction.
\end{proof}

\begin{lemma}[Cancellation of a nonempty identity factor]
    \label{lem:matrix_kronecker_identity_cancellation}
    \lean{Matrix.kronecker_one_injective, Matrix.mul_kronecker_one,
        Matrix.mul_eq_one_of_kronecker_one,
        Matrix.mul_conjTranspose_eq_one_of_kronecker_one,
        Matrix.conjTranspose_mul_eq_one_of_kronecker_one}
    \leanok
    Let $D>0$. The map $A\mapsto A\otimes I_D$ is injective and preserves
    matrix products. Consequently,
    $(F\otimes I_D)(G\otimes I_D)=I$ implies $FG=I$.
    In particular, orthonormal rows or columns of $F\otimes I_D$ imply the
    corresponding orthonormality of $F$.
\end{lemma}
\begin{proof}\leanok
    Evaluate the amplified identity on one basis vector of the nonzero
    factor. Multiplicativity of the Kronecker product gives the product and
    adjoint specializations.
\end{proof}

\begin{definition}[Fusion-isometry family]%
    \label{def:mpdo_bnt_fusion_isometry_family}
    \lean{MPOTensor.BNTFusionIsometryFamily}
    \leanok
    \uses{def:mpdo_operator_product, def:mpdo_diagonal_chi_family}
    A \emph{fusion-isometry family} over a finite set of labels consists of a
    family of tensors $M_\gamma$ with a common physical dimension and
    per-label bond dimensions, diagonal matrices $\chi_{\alpha,\beta,\gamma}$
    with positive entries, and isometries $U_{\alpha,\beta}$ such that, site
    by site,
    \begin{align}
        U_{\alpha,\beta}(M_\alpha M_\beta)^{ij}
            U_{\alpha,\beta}^\dagger
        &=
        \bigoplus_\gamma
            \chi_{\alpha,\beta,\gamma}\otimes M_\gamma^{ij}.
        \label{eq:rfp_fusion_identity}
    \end{align}
    This is the isometry statement of
    \cite[Theorem~4.14(iii)]{Cirac2016MPDO_arXiv}; the tensors are the
    vertically read basis of normal tensors of
    \cite[Proposition~4.13]{Cirac2016MPDO_arXiv}, so the physical dimension
    is the bond dimension of the original tensor. Statement (iii) of the
    source theorem additionally asserts the idempotent identity for the trace
    scalars of the vertical decomposition, recorded separately by
    Definition~\ref{def:mpdo_bnt_label_idempotent_coefficient_form}.

    \textbf{Scope restriction (full-support fusion family):} The displayed
    family assumes $U_{\alpha,\beta}^\dagger U_{\alpha,\beta}=1$ on the whole
    product bond space. The unrestricted source statement permits a common
    zero corner and instead gives a coisometry onto the active direct sum,
    together with exact reconstruction. This distinction is documented in
    \path{docs/paper-gaps/cpsv16_figure11_fusion_coisometry.tex}.
\end{definition}

\begin{definition}[Active-support fusion coisometry family]
    \label{def:mpdo_bnt_fusion_coisometry_family}
    \lean{MPOTensor.BNTFusionCoisometryFamily}
    \lean{MPOTensor.BNTFusionCoisometryFamily.%
        toBNTFusionIsometryFamily_of_fullSupport}
    \leanok
    \uses{def:mpdo_bnt_fusion_isometry_family}
    An \emph{active-support fusion coisometry family} consists of the same
    labelled tensors and positive diagonal matrices, together with maps
    $U_{\alpha,\beta}$ satisfying
    $U_{\alpha,\beta}U_{\alpha,\beta}^\dagger=1$, the forward fusion identity
    \begin{align}
        U_{\alpha,\beta}(M_\alpha M_\beta)^{ij}
          U_{\alpha,\beta}^\dagger
        &=G_{\alpha,\beta}^{ij},
        \notag
    \end{align}
    and exact reconstruction
    \begin{align}
        (M_\alpha M_\beta)^{ij}
        &=U_{\alpha,\beta}^\dagger
          G_{\alpha,\beta}^{ij}U_{\alpha,\beta}.
        \label{eq:rfp_coisometry_reconstruction}
    \end{align}
    Zero-dimensional retained sums are permitted. If in addition
    $U_{\alpha,\beta}^\dagger U_{\alpha,\beta}=1$ for every pair, this family
    gives the preceding full-support fusion-isometry family.

    \textbf{Local fix (Figure-11 fusion coisometry):} The source uses the
    retained-row orientation of Proposition~4.13. Thus its fusion maps are
    coisometries onto the active sectors; exact reconstruction records the
    discarded common zero corner. This convention is documented in
    \path{docs/paper-gaps/cpsv16_figure11_fusion_coisometry.tex}.
\end{definition}

For fixed labels $\alpha,\beta$, put
$G_{\alpha,\beta}^{ij}=\bigoplus_\gamma
\chi_{\alpha,\beta,\gamma}\otimes M_\gamma^{ij}$. Then
\cite[Theorem~4.14(iii), equation~(89), lines~986--991]
{Cirac2016MPDO_arXiv} gives
% For fixed \alpha,\beta,
%   U_{\alpha,\beta}(M_\alpha M_\beta)^{ij}U_{\alpha,\beta}^{\dagger}
%     = G_{\alpha,\beta}^{ij},
%   G_{\alpha,\beta}^{ij}
%     = \bigoplus_\gamma \chi_{\alpha,\beta,\gamma}\otimes M_\gamma^{ij}.
% Coefficientwise, with x,y,i,j free,
%   (G_{\alpha,\beta}^{ij})_{x,y}
%     = \sum_{q,a_1,a_2,b_1,b_2}
%       (U_{\alpha,\beta})_{x,(a_1,b_1)}
%       (M_\alpha^{iq})_{a_1,a_2}(M_\beta^{qj})_{b_1,b_2}
%       (U_{\alpha,\beta}^{\dagger})_{(a_2,b_2),y}.
% Thus q,a_1,a_2,b_1,b_2 are contracted.  For
% x=(\gamma,r,a) and y=(\gamma',r',a'),
%   (G_{\alpha,\beta}^{ij})_{(\gamma,r,a),(\gamma',r',a')}
%     = \delta_{\gamma,\gamma'}(\chi_{\alpha,\beta,\gamma})_{r,r'}
%       (M_\gamma^{ij})_{a,a'},
% or, since chi is diagonal,
%   \delta_{\gamma,\gamma'}\delta_{r,r'}\chi_{\alpha,\beta,\gamma;r}
%       (M_\gamma^{ij})_{a,a'}.
% The sector \gamma is a direct-sum block label, not a contraction index.
% Ink-to-index map: the combined west and east fuser legs are x and y;
% the north and south physical legs are i and j; the unique inter-row
% pairleg is q; and the four fuser-to-tensor strands carry a_1,b_1,a_2,b_2.
% The left and right objects each have boundary signature (1,1,1,1).
% The single G_{\alpha,\beta} atom has no internal contraction.  The cited
% source states the algebraic identity and direct-sum decomposition, but
% contains no dedicated diagram for U_{\alpha,\beta}.
\par\noindent\hfil
\tnpic[rows={op,op}]{
    \tnfuse[combined=west]{U_{\alpha,\beta}} &
    \tn[mpo, up=$i$]{M_\alpha} &
    \tnfuse[combined=east]{U_{\alpha,\beta}^{\dagger}} \\
    & \tn[mpo, down=$j$]{M_\beta} &
}
\qquad$=$\qquad
\tnpic[physical=updown]{
    \tn[mpo, up=$i$, down=$j$]{G_{\alpha,\beta}}
}
\hfil\par

\begin{lemma}[Weighted zipper identities on the active support]
    \label{lem:mpdo_bnt_active_support_weighted_zipper}
    \lean{MPOTensor.BNTFusionCoisometryFamily.%
        fusionCoisometry_mul_mulTensor}
    \lean{MPOTensor.BNTFusionCoisometryFamily.%
        mulTensor_mul_fusionCoisometry_conjTranspose}
    \leanok
    \uses{def:mpdo_bnt_fusion_coisometry_family}
    An active-support fusion coisometry satisfies
    \begin{align}
        U_{\alpha,\beta}(M_\alpha M_\beta)^{ij}
        &=G_{\alpha,\beta}^{ij}U_{\alpha,\beta},
        \notag\\
        (M_\alpha M_\beta)^{ij}U_{\alpha,\beta}^\dagger
        &=U_{\alpha,\beta}^\dagger G_{\alpha,\beta}^{ij}.
        \notag
    \end{align}
\end{lemma}
\begin{proof}\leanok
    Substitute exact reconstruction
    (\ref{eq:rfp_coisometry_reconstruction}) into the left-hand sides and use
    $U_{\alpha,\beta}U_{\alpha,\beta}^\dagger=1$.
\end{proof}

\begin{theorem}[Weighted zipper identities]%
    \label{thm:mpdo_bnt_weighted_zipper}
    \lean{MPOTensor.BNTFusionIsometryFamily.fusionIsometry_mul_mulTensor}
    \lean{MPOTensor.BNTFusionIsometryFamily.%
        mulTensor_mul_fusionIsometry_conjTranspose}
    \leanok
    \uses{def:mpdo_bnt_fusion_isometry_family}
    For every pair of labels and physical indices, write
    $G_{\alpha,\beta}^{ij}
    =\bigoplus_\gamma\chi_{\alpha,\beta,\gamma}
    \otimes M_\gamma^{ij}$.
    The fusion isometry satisfies the two chi-weighted zipper identities
    \begin{align}
        U_{\alpha,\beta}(M_\alpha M_\beta)^{ij}
        &=G_{\alpha,\beta}^{ij}U_{\alpha,\beta},
        \notag\\
        (M_\alpha M_\beta)^{ij}U_{\alpha,\beta}^\dagger
        &=U_{\alpha,\beta}^\dagger G_{\alpha,\beta}^{ij}.
        \notag
    \end{align}
    The corresponding identities in
    \cite[Equation~(21)]{Bultinck2017Anyons} have identity matrices on
    the multiplicity spaces in place of the positive diagonal matrices
    $\chi$. Their unweighted form requires the length-independent integer
    specialization and is not a consequence of the present hypotheses. Even
    after that specialization, comparison of the two triple-fusion orders and
    final-label separation are still needed to construct an invertible
    $F$-move.
\end{theorem}
\begin{proof}\leanok
    Multiply the fusion identity (\ref{eq:rfp_fusion_identity}) on the right
    by $U_{\alpha,\beta}$ and on the left by
    $U_{\alpha,\beta}^\dagger$, respectively, and use
    $U_{\alpha,\beta}^\dagger U_{\alpha,\beta}=1$.
\end{proof}

\begin{theorem}[Length-independent zipper identities and completeness]%
    \label{thm:mpdo_bnt_length_independent_zipper}
    \lean{MPOTensor.BNTFusionIsometryFamily.%
        chi_matrix_eq_one_of_lengthIndependent}
    \lean{MPOTensor.BNTFusionIsometryFamily.%
        fusionIsometry_mul_mulTensor_of_lengthIndependent}
    \lean{MPOTensor.BNTFusionIsometryFamily.%
        mulTensor_mul_fusionIsometry_conjTranspose_of_lengthIndependent}
    \lean{MPOTensor.BNTFusionIsometryFamily.%
        mulTensor_eq_conjTranspose_mul_unweightedDirectSum_mul_of_lengthIndependent}
    \leanok
    \uses{def:mpdo_bnt_fusion_isometry_family,
        def:mpdo_bnt_label_length_independent,
        thm:mpdo_bnt_label_length_independent_natural,
        thm:mpdo_bnt_weighted_zipper}
    Suppose the structure coefficients have positive trace-power form and are
    independent of the positive chain length. Then
    $\chi_{\alpha,\beta,\gamma}=1_{r_{\alpha,\beta,\gamma}}$ for every triple
    of labels. If
    \begin{align}
        H_{\alpha,\beta}^{ij}
        &=
        \bigoplus_\gamma
            1_{r_{\alpha,\beta,\gamma}}\otimes M_\gamma^{ij},
        \notag
    \end{align}
    the fusion isometries satisfy
    \begin{align}
        U_{\alpha,\beta}(M_\alpha M_\beta)^{ij}
        &=H_{\alpha,\beta}^{ij}U_{\alpha,\beta},
        \notag\\
        (M_\alpha M_\beta)^{ij}U_{\alpha,\beta}^\dagger
        &=U_{\alpha,\beta}^\dagger H_{\alpha,\beta}^{ij},
        \notag\\
        (M_\alpha M_\beta)^{ij}
        &=U_{\alpha,\beta}^\dagger
            H_{\alpha,\beta}^{ij}U_{\alpha,\beta}.
        \label{eq:rfp_unweighted_reconstruction}
    \end{align}
    These are the identity-weight specialization of Equations~(19)--(21) in
    \cite[Section~3]{Bultinck2017Anyons}. The MPDO fusion identity appears in
    \cite[Theorem~4.14(iii), lines~986--991]{Cirac2016MPDO_arXiv}; its
    identity-weight specialization follows from the length-independent case
    at \cite[line~1010]{Cirac2016MPDO_arXiv}.
\end{theorem}
\begin{proof}\leanok
    Length independence makes every positive diagonal entry of
    $\chi_{\alpha,\beta,\gamma}$ equal to one. Substitute the resulting
    identity matrices into
    Theorem~\ref{thm:mpdo_bnt_weighted_zipper}. For completeness, insert
    $U_{\alpha,\beta}^\dagger U_{\alpha,\beta}=1$ on both sides of the product
    letter and apply the fusion identity between the two inserted factors,
    yielding (\ref{eq:rfp_unweighted_reconstruction}).
\end{proof}

\begin{theorem}[Common positive blocking of the BNT fusion isometries]
    \label{thm:mpdo_bnt_common_positive_blocked_fusion}
    \lean{MPSTensor.wordTupleSpanTop_reindexPhysical_equiv}
    \lean{MPSTensor.WordTupleSpanTop.isInjective_one}
    \lean{MPOTensor.BNTFusionIsometryFamily.%
        blockedUnweightedDirectSumLetter}
    \lean{MPOTensor.BNTFusionIsometryFamily.%
        blocked_fusion_of_lengthIndependent}
    \lean{MPOTensor.BNTFusionIsometryFamily.%
        fusionIsometry_mul_blocked_mulTensor_of_lengthIndependent}
    \lean{MPOTensor.BNTFusionIsometryFamily.%
        blocked_mulTensor_mul_fusionIsometry_conjTranspose_of_lengthIndependent}
    \lean{MPOTensor.BNTFusionIsometryFamily.%
        blocked_mulTensor_eq_conjTranspose_mul_unweightedDirectSum_mul}
    \lean{MPOTensor.BNTFusionIsometryFamily.%
        exists_positive_block_with_injective_fusion}
    \leanok
    \uses{thm:cpsv_bnt_blocked_simultaneous_span,
        thm:mpdo_positive_blocking_doubled_index,
        thm:mpdo_positive_blocking_product,
        thm:mpdo_bnt_length_independent_zipper}
    Suppose that the labelled doubled-index tensors form the basis of normal
    tensors, and that their fusion-isometry family has positive trace-power
    coefficients independent of the positive chain length. Then there is
    one integer $L>0$ such that the blocked tensors $M_\gamma^{[L]}$ span
    their full product matrix algebra in one letter. Hence every
    $M_\gamma^{[L]}$ is injective.

    For this same $L$, write
    \begin{align}
        P_{\alpha,\beta}^{IJ}
        &=(M_\alpha^{[L]}M_\beta^{[L]})^{IJ},
        \notag\\
        H_{\alpha,\beta}^{IJ}
        &=\bigoplus_\gamma
          1_{r_{\alpha,\beta,\gamma}}
            \otimes(M_\gamma^{[L]})^{IJ},
        \notag
    \end{align}
    where $r_{\alpha,\beta,\gamma}$ is the dimension of the specialized
    chi space. The original fusion isometry satisfies
    \begin{align}
        U_{\alpha,\beta}P_{\alpha,\beta}^{IJ}
            U_{\alpha,\beta}^\dagger
        &=H_{\alpha,\beta}^{IJ},
        \label{eq:rfp_blocked_fusion}
    \end{align}
    the two blocked zipper identities
    \begin{align}
        U_{\alpha,\beta}P_{\alpha,\beta}^{IJ}
        &=H_{\alpha,\beta}^{IJ}U_{\alpha,\beta},
        \notag\\
        P_{\alpha,\beta}^{IJ}U_{\alpha,\beta}^\dagger
        &=U_{\alpha,\beta}^\dagger H_{\alpha,\beta}^{IJ},
        \notag
    \end{align}
    and the literal reconstruction
    \begin{align}
        P_{\alpha,\beta}^{IJ}
        &=U_{\alpha,\beta}^\dagger
          H_{\alpha,\beta}^{IJ}U_{\alpha,\beta}.
        \notag
    \end{align}
\end{theorem}

\begin{proof}\leanok
    Choose the positive block length supplied by simultaneous BNT blocking.
    The canonical pairing of the blocked ket and bra words transfers the
    one-letter product-algebra span to the blocked MPO tensors, and projection
    to each labelled factor gives injectivity. Along a nonempty block, the
    products of the conjugated fusion letters telescope using only
    $U_{\alpha,\beta}^\dagger U_{\alpha,\beta}=1$. Length independence
    replaces every chi matrix by the identity on its own chi space. The two
    zipper identities and reconstruction then follow by multiplying the
    blocked fusion identity (\ref{eq:rfp_blocked_fusion}) by
    $U_{\alpha,\beta}$ or $U_{\alpha,\beta}^\dagger$ on the appropriate side.
\end{proof}

The three rows below are, respectively, the two unweighted zipper identities
and their local reconstruction consequence. The symbol
$H$ denotes $H_{\alpha,\beta}^{ij}$, and the open physical legs carry the
indices $i,j$. The reconstruction uses only
$U_{\alpha,\beta}^{\dagger}U_{\alpha,\beta}=1$; it does not assert
$U_{\alpha,\beta}U_{\alpha,\beta}^{\dagger}=1$, surjectivity, or completeness
of the total fusion space. The zipper orientation is that of
\cite[Section~3, Equations~(19)--(21)]{Bultinck2017Anyons}; the MPDO identity
is \cite[Theorem~4.14(iii), lines~986--991]{Cirac2016MPDO_arXiv}, with identity
weights in the length-independent case at
\cite[line~1010]{Cirac2016MPDO_arXiv}.
% Coefficient formula: \(P_{\alpha,\beta}^{ij}=(M_\alpha M_\beta)^{ij}
%   =\sum_q M_\alpha^{iq}\otimes M_\beta^{qj}\).  The middle physical index q
%   is contracted once; i and j are the one free north and south indices.
% The three rows are \(UP=HU\), \(PU^\dagger=U^\dagger H\), and
%   \(P=U^\dagger HU\).  The final equality cancels \(U^\dagger U\) only;
%   it does not assert \(UU^\dagger=1\) or completeness of the total fusion
%   space.
% Boundary signatures (west,east,up,down), left and right respectively:
%   row 1: (1,2,1,1)=(1,2,1,1);
%   row 2: (2,1,1,1)=(2,1,1,1);
%   row 3: (2,2,1,1)=(2,2,1,1).
% Ink-to-index map: each M box carries one product factor; their vertical
%   stroke is the q contraction and their outer legs are i,j.  Each H box is
%   the single aggregate letter \(H_{\alpha,\beta}^{ij}\).  A fusion bar
%   presents the separate \(\alpha,\beta\) ports at \(\{1,2\}\) and the fused
%   port at center.  The
%   H--bar bonds lie on a fused row; the second row is structural, and the
%   reconstruction hole prevents it from adding a second H channel.
% Source verdict: Bultinck et al. arXiv:1511.08090, Eq. (21), states the two
%   zipper orientations algebraically but gives no dedicated figure for that
%   equation.  CPSV16 Theorem 4.14(iii) supplies the MPDO identity.  The old
%   catalogue picture fixes only row order and topology, not coordinates;
%   therefore the current diagonal fused bonds preserve the asserted
%   contraction without inventing a centred tall H atom.
\par\noindent\hfil
\tnpic[rows={op,op}]{
  \tnfuse[span=2, west at=center, east at={1,2}]{U_{\alpha,\beta}} &
  \tn[mpo, up=$i$]{M_\alpha} \\
  & \tn[mpo, down=$j$]{M_\beta}
}
\qquad$=$\qquad
\tnpic[rows={op:fused,wire:west none}]{
  \tn[mpo, west at=center, east at=center, up=$i$, down=$j$]{H} &
  \tnfuse[span=2, west at=center, east at={1,2}]{U_{\alpha,\beta}} \\
  &
}
\hfil\par
\par\noindent\hfil
\tnpic[rows={op,op}]{
  \tn[mpo, up=$i$]{M_\alpha} &
  \tnfuse[span=2, west at={1,2}, east at=center]
    {U_{\alpha,\beta}^{\dagger}} \\
  \tn[mpo, down=$j$]{M_\beta} &
}
\qquad$=$\qquad
\tnpic[rows={op:fused,wire:east none}]{
  \tnfuse[span=2, west at={1,2}, east at=center]
    {U_{\alpha,\beta}^{\dagger}} &
  \tn[mpo, west at=center, east at=center, up=$i$, down=$j$]{H} \\
  &
}
\hfil\par
\par\noindent\hfil
\tnpic[rows={op,op}]{
  \tn[mpo, up=$i$]{M_\alpha} \\
  \tn[mpo, down=$j$]{M_\beta}
}
\qquad$=$\qquad
\tnpic[rows={op:fused,wire}]{
  \tnfuse[span=2, west at={1,2}, east at=center]
    {U_{\alpha,\beta}^{\dagger}} &
  \tn[mpo, west at=center, east at=center, up=$i$, down=$j$]{H} &
  \tnfuse[span=2, west at=center, east at={1,2}]{U_{\alpha,\beta}} \\
  & \tnskip &
}
\hfil\par
